\scanpage{170}
De plus

a) $\varphi_*$ commute aux p.d. ssi $B'_i=A'_1{}^{(i-1)}B'_1\quad\forall i\geq1$.

b) $\varphi_*$ commute aux graduations extérieures ssi $B'_i=0\quad\forall i\geq1$.

c) $\varphi_*$ commute aux filtrations croissantes grossières
(dont les termes sont des sous-$S$ ou sous-$S'$-modules) ssi
$A'_1=\lambda'X_0\quad(\lambda'\in k'\{T\})$ et
\[
 B'_i\unsure{=}\sum_{0\leq j\leq i-1}b'_{ij}X_0^{(j)\deleted{\unsure{*}}},
 \qquad b'_{ij}\in k'\{T\}\quad\text{(pour mémoire)}
\]

d) $\varphi_*$ commute aux filtrations \deleted{\unsure{déc.}} croissantes
fines [dont les termes sont des sous-$k$ (ou $k'$)
modules, \unsure{-- les sous|non des sous} $\unsure{S}$ (ou $S'$)-modules
\[
 \operatorname{Fil}^{\mathrm{gross}}_{\unsure{d}} \mathcal A_{k,n}
 =\left\{\sum F_{i_0\cdots i_r}\,dX_{i_0}\wedge\cdots\wedge dX_{i_r}
 \;\middle|\;F_{i_0\cdots i_r}\in k\{X_0,\ldots,X_n\}
 \text{ de ``degré'' }\leq d-r\right\}
\]
\margin{\uline{NB} La filtration grossière est la filtration par
$S$-sous-modules engendrée par la filtration fine}

\unsure{donc} c'est la filtration croissante la plus fine pour laquelle
\deleted{\unsure{k}} $k\subset\operatorname{Fil}_0$, $X_r,dX_r\in\operatorname{Fil}_1$,
\[
 \operatorname{Fil}_i\cdot\operatorname{Fil}_j\subset\operatorname{Fil}_{i+j},
 \quad\text{et}\quad x\in\operatorname{Fil}_d\Longrightarrow x^{(n)}\in\operatorname{Fil}_{dn};
\]
c'est aussi la filtration croissante associée à la graduation par le degré total de
\[
 (\mathcal A_k)_n\simeq k\{X_0,\ldots,X_n\}[dX_0,\ldots,dX_n]/\langle\textstyle\sum dX_i\rangle
\]
(degré compatible aux p.d. et tel que $X_r$ et $dX_r$ de degré $1$ -- donc
c'est aussi le degré pour lequel $(\mathcal A_k)_n$ devient de façon naturelle un
module gradué sur l'anneau gradué $S$, avec différentielle de degré $0$.)]
\deleted{il f. et \inserted{ssi} qu'on ait}
\[
 \left\{\begin{aligned}
 &\varphi_0(k)\subset k',\\
 &A'_1=\unsure{\lambda'_0|\alpha'_0}X_0\qquad\unsure{\lambda'_0|\alpha'_0}\in k',\\
 &B'_i=\sum_{\substack{j,h\geq0\\j+h\leq i-1}}\alpha'_{i,jh}\,X_0^{(j)}X_1^{(h)}
 \end{aligned}\right.
 \qquad\figure{\begin{tikzpicture}\node[draw,ellipse] {$(\unsure{\forall i\geq1})$};\end{tikzpicture}}
\]
\margin{(Cette condition est \unsure{vide|vraie} si $2$ inv. dans $k'$, car les $B'_i$ sont nuls!)}
\margin{169}
